% Part: lambda-calculus
% Chapter: lambda-definability
% Section: fixpoints

\documentclass[../../../include/open-logic-section]{subfiles}

\begin{document}

\olfileid{lam}{ldf}{fp}

\olsection{স্থির-বিন্দুসমূহ}

ধরা যাক, আমরা পুনরাবৃত্তির সাহায্যে ক্রমগুণিতক অপেক্ষকটিকে এমন একটি
পদ~$\fn{Fac}$ হিসেবে সংজ্ঞায়িত করতে চাই, যার নিচের ধর্মটি আছে:
\[
\fn{Fac} \ident \lambd[n][\fn{IsZero}\, n\, \num 1 (\fn{Mult}\, n (\fn{Fac}(\fn{Pred} \, n)))]
\]
অর্থাৎ $n = 0$ হলে $n$-এর ক্রমগুণিতক $1$, আর অন্যথায় তা $n$ গুণ
$n-1$-এর ক্রমগুণিতক। অবশ্যই এভাবে $\fn{Fac}$ পদটি সংজ্ঞায়িত করা যায়
না, কারণ ডান পাশেই $\fn{Fac}$ আছে। আত্মউল্লেখসংবলিত এমন পুনরাবৃত্ত
সংজ্ঞা ল্যাম্বডা ক্যালকুলাসের অংশ নয়। যেমন, কোনো পদকে নিচের মতো
সংজ্ঞায়িত করলে
\[
\fn{Mult} \ident \lambd[ab][a (\fn{Add}\, b) \num{0}]
\]
ডান পাশে শুধু $\fn{Add}$-এর মতো আগে সংজ্ঞায়িত পদ থাকে। $\fn{Add}$-এর
জায়গায় তার সংজ্ঞাকারী পদ বসিয়ে আমরা সবসময় সেটিকে বাদ দিতে পারি। এতে
$\fn{Mult}$ একটি বিশুদ্ধ ল্যাম্বডা পদ হিসেবে পাওয়া যাবে; $\fn{Add}$
নিজেও যদি কোনো সংজ্ঞায়িত পদ ব্যবহার করত (যেমন $\fn{Add}'$ করে), তবে
প্রক্রিয়াটি চালিয়ে শেষ পর্যন্ত একটি বিশুদ্ধ ল্যাম্বডা পদে পৌঁছানো যেত।
% BN-SRC-311 TARGET-CORRECTION-BEGIN
\par\smallskip\noindent\textnormal{\small\textbf{উৎস-সংশোধন (BN-SRC-311):} হিমায়িত উদাহরণটি $\lambd[ab][a(\fn{Add}\,a)0]$ লেখে; এতে আবদ্ধ~$b$ অব্যবহৃত থেকে যায়, ফলে পদটি গুণফলের বদলে~$a$-এর বর্গ গণনা করে, এবং কাঁচা~$0$ বিশুদ্ধ ল্যাম্বডা পদও নয়। ঠিক আগের পাটীগাণিতিক অংশের সংশোধিত নির্মাণ অনুসারে বাংলা উদাহরণে $\lambd[ab][a(\fn{Add}\,b)\num{0}]$ পুনরুদ্ধার করা হয়েছে।} % BN-SRC-311 TARGET-CORRECTION-END

কিন্তু উপরের $\fn{Fac}$-এর মতো পুনরাবৃত্ত সংজ্ঞার ক্ষেত্রে এটি সত্য নয়।
ডান পাশের $\fn{Fac}$-এর জায়গায় $\fn{Fac}$-এর সংজ্ঞাটিই বসালে পাই:
\begin{align*}
  \fn{Fac} & \ident \lambd[n][\fn{IsZero}\, n\, \num{1}] \\
    & \qquad (\fn{Mult}\, n
  ((\lambd[n][\fn{IsZero} \, n \, \num 1\, (\fn{Mult}\, n\, (\fn{Fac}
    (\fn{Pred}\, n)))]) (\fn{Pred}\, n)))
\end{align*}
তবুও ডান পাশ থেকে $\fn{Fac}$ বাদ যায়নি। স্পষ্টতই, প্রক্রিয়াটি আবার
করলে সংজ্ঞাটি ক্রমশ দীর্ঘ হবে এবং কখনোই কোনো বিশুদ্ধ ল্যাম্বডা পদ দেবে
না। তাই ক্রমগুণিতককে (বা আরও সাধারণভাবে পুনরাবৃত্ত অপেক্ষককে) এভাবে
সংজ্ঞায়িত করা কার্যকর নয়।

তবে পুনরাবৃত্ত সংজ্ঞাটি আমাদের কিছু তথ্য দেয়। $f$ যদি ক্রমগুণিতক
অপেক্ষকের উপস্থাপক কোনো পদ হতো, তাহলে
\[
\fn{Fac}' \ident \lambd[g][\lambd[n][\fn{IsZero} \, n \, \num 1\, (\fn{Mult} \, n\, (g (\fn{Pred} n)))]]
\]
পদটিকে $f$-এর উপর প্রয়োগ করলে, অর্থাৎ $\fn{Fac}'\,f$, সেটিও
ক্রমগুণিতক অপেক্ষকের উপস্থাপক হতো। অর্থাৎ $\fn{Fac}'$-কে একটি অপেক্ষক
গ্রহণ করে আরেকটি অপেক্ষক ফেরত দেওয়া অপেক্ষক হিসেবে দেখলে,
$f$ ক্রমগুণিতকের উপস্থাপক হওয়ার শর্তে $\fn{Fac'}\, f$-এর মান ঠিক~$f$।
$\fn{Fac}' \, f \equal[\beta] f$ ধর্মের কোনো অপেক্ষক~$f$-কে
$\fn{Fac}'$-এর একটি \emph{স্থির-বিন্দু} বলা হয়। অতএব ক্রমগুণিতক
$\fn{Fac}'$-এর একটি স্থির-বিন্দু।

ল্যাম্বডা ক্যালকুলাসে এমন পদ আছে, যা কোনো প্রদত্ত পদের স্থির-বিন্দু
গণনা করে; সেগুলি ব্যবহার করে $\fn{Fac}'$-এর মতো একটি পদকে ক্রমগুণিতকের
সংজ্ঞায় পরিণত করা যায়।

\begin{defn}
  \ollabel{defn:Turing-Y}
  \emph{Y-সমাবেশক} হলো নিচের পদটি:
  \[
  Y \ident (\lambd[ux][x(uux)])(\lambd[ux][x(uux)]).
  \]
\end{defn}

\begin{thm}
  যেকোনো পদ~$g$-এর জন্য $Yg \red g(Yg)$---এই ধর্মটি $Y$-এর আছে। অতএব
  $Yg$ সর্বদা~$g$-এর একটি স্থির-বিন্দু।
\end{thm}

\begin{proof}
  $(\lambd[ux][x(uux)])$-কে সংক্ষেপে~$U$ লিখি, যাতে $Y \ident UU$। তাহলে
  \begin{align*}
    Y g &\ident (\lambd[ux][x(uux)])U\,g \\
    &\red (\lambd[x][x(UUx)])g\\
    & \red g(UUg) \ident g(Yg).
  \end{align*}
  যেহেতু $g(Yg)$ ও~$Yg$ উভয়ই $g(Yg)$-তে হ্রাস পায়, তাই $g(Yg)
  \equal[\beta] Yg$; অতএব $Yg$ হলো~$g$-এর একটি স্থির-বিন্দু।
\end{proof}

অবশ্য $Yg$ একটি রিডেক্স বলে হ্রাস অনির্দিষ্টকাল চলতে পারে:
\begin{align*}
  Y g &\red g (Y g) \\
    &\red g (g (Y g)) \\
    &\red g(g (g (Y g)))\\
    &\ldots
\end{align*}
তাই $Yg$-কে নিজের উপর~$g$-এর অসীম পুনরাবৃত্ত প্রয়োগ হিসেবে ভাবা যায়।
এর উপর $g$-কে আরও একবার প্রয়োগ করলে---কথার কথা হিসেবে---নতুন কিছুই ঘটে
না: $Yg$-এর উপর অসীমবার~$g$ প্রয়োগ করে পাওয়া পদের উপর আরেকবার~$g$
প্রয়োগ করলেও সেটি $Yg$-এর উপর অসীমবার~$g$ প্রয়োগের ফলই থাকে।

লক্ষ করো, $Yg$ থেকে শুরু হওয়া উপরের $\beta$-হ্রাসের অনুক্রমটি অসীম।
তাই $Yg$-কে কোনো পদে প্রয়োগ করে $(Yg)N$ নিলে, সেটিও অসীমসংখ্যক ভিন্ন
পদে হ্রাস পাবে: $(g(Yg))N$, $(g(g(Yg)))N$, \dots। তা সত্ত্বেও হ্রাসের
\emph{অন্য} কোনো অনুক্রম স্বাভাবিক রূপে গিয়ে শেষ হতে পারে।

উদাহরণ হিসেবে ক্রমগুণিতকই নাও। $\fn{Fac}$-কে $Y\, \fn{Fac}'$ (অর্থাৎ
$\fn{Fac'}$-এর একটি স্থির-বিন্দু) হিসেবে সংজ্ঞায়িত করো। তাহলে:
\begin{align*}
  \fn{Fac}\,\num 3
  &\red Y\, \fn{Fac}' \, \num 3 \\
  &\red \fn{Fac}' (Y \,\fn{Fac}') \, \num 3 \\
  &\ident (\lambd[x][\lambd[n][\fn{IsZero} \, n\, \num 1\,
      (\fn{Mult}\, n\, (x (\fn{Pred}\, n)))]]) \, \fn{Fac} \, \num 3\\
  & \red \fn{IsZero} \, \num 3\, \num 1\,
      (\fn{Mult}\, \num 3\, (\fn{Fac} (\fn{Pred}\, \num 3))) \\
  &\red  \fn{Mult}\, \num 3 \, (\fn{Fac} \, \num 2).
  \intertext{একইভাবে,}
  \fn{Fac}\,\num 2
  &\red  \fn{Mult}\, \num 2 \, (\fn{Fac} \, \num 1) \\
  \fn{Fac}\,\num 1
  &\red  \fn{Mult}\, \num 1 \, (\fn{Fac} \, \num 0)
  \intertext{কিন্তু}
  \fn{Fac}\,\num 0
  &\red \fn{Fac}' (Y \,\fn{Fac}') \, \num 0 \\
  &\ident (\lambd[x][\lambd[n][\fn{IsZero} \, n\, \num 1\,
      (\fn{Mult}\, n\, (x (\fn{Pred}\, n)))]]) \, \fn{Fac} \, \num 0\\
  & \red \fn{IsZero} \, \num 0\, \num 1\,
      (\fn{Mult}\, \num 0\, (\fn{Fac} (\fn{Pred}\, \num 0))).\\
  &\red \num 1.
  \intertext{অতএব সব মিলিয়ে}
  \fn{Fac}\,\num 3
  &\red  \fn{Mult}\, \num 3 \,
  (\fn{Mult} \, \num 2\, (\fn{Mult} \, \num 1\, \num 1)).
\end{align*}

$\fn{Fac'}$-এর ক্ষেত্রে যা ঘটে, যেকোনো পুনরাবৃত্ত সংজ্ঞার ক্ষেত্রেও তা-ই।
ধরা যাক, আমাদের একটি পুনরাবৃত্ত সমীকরণ আছে
\begin{align*}
g\,x_1\dots x_n & \equal[\beta] N
\intertext{যেখানে $N$-এর মধ্যে~$g$ ও~$x_1$, \dots,~$x_n$ থাকতে পারে। তাহলে
এমন একটি পদ~$G \ident (Y \lambd[g][\lambd[x_1\dots x_n][N]])$ সর্বদাই
আছে, যাতে}
G\,x_1 \dots x_n & \equal[\beta] \Subst{N}{G}{g}.
\intertext{কারণ স্থির-বিন্দু উপপাদ্য অনুসারে,}
G \ident (Y \lambd[g][\lambd[x_1\dots x_n][N]]) & \red \lambd[g][\lambd[x_1\dots x_n][N]](Y \lambd[g][\lambd[x_1\dots x_n][N]])\\
& \ident (\lambd[g][\lambd[x_1\dots x_n][N]])\,G
\intertext{এবং তাই}
G\,x_1 \dots x_n
& \red (\lambd[g][\lambd[x_1\dots x_n][N]])\,G\,x_1\dots x_n\\
& \red (\lambd[x_1\dots x_n][\Subst{N}{G}{g}])\,x_1\dots x_n\\
  & \red \Subst{N}{G}{g}.
\end{align*}

\olref{defn:Turing-Y}-এর $Y$-সমাবেশকটি অ্যালান টুরিংয়ের। অ্যালোন্‌জো
চার্চ আরেকটি রূপ প্রস্তাব করেছিলেন; সেটিকে~$Y_C$ বলব:
\[
Y_C \ident \lambd[g][(\lambd[x][g(xx)])(\lambd[x][g(xx)])].
\]
চার্চের সমাবেশকটি টুরিংয়েরটির চেয়ে একটু দুর্বল: $Y_C g
\equal[\beta] g(Y_C g)$, কিন্তু $Y_C g \bred g(Y_C g)$ নয়। $V$-কে
$\lambd[x][g(xx)]$ পদটি ধরো, যাতে $Y_C \ident \lambd[g][VV]$। তাহলে
\begin{align*}
  VV & \ident (\lambd[x][g(xx)])V \red g(VV)
  \text{ এবং তাই}\\
  Y_C g & \ident (\lambd[g][VV])g \red VV \red g(VV), \text{ কিন্তু আবার}\\
  g(Y_C g) & \ident g((\lambd[g][VV])g) \red g(VV).
\end{align*}
অন্যভাবে বললে, $Y_Cg$ ও~$g(Y_Cg)$ একই পদ~$g(VV)$-তে হ্রাস পায়; তাই
$Y_Cg \equal[\beta] g(Y_Cg)$। প্রয়োগের জন্য এটিই প্রায়ই যথেষ্ট।
% BN-SRC-312 TARGET-CORRECTION-BEGIN
\par\smallskip\noindent\textnormal{\small\textbf{উৎস-সংশোধন (BN-SRC-312):} হিমায়িত তুলনামূলক বাক্যটি চার্চের~$Y_C$ নিয়ে কথা বলেও দুইবার টুরিংয়ের~$Y$ লিখেছে এবং তাই ঠিক আগের উপপাদ্যের প্রত্যক্ষ $Yg\bred g(Yg)$ ধর্মকেই অস্বীকার করেছে। বাংলা বাক্যে উভয় স্থানে~$Y_C$ পুনরুদ্ধার করা হয়েছে; নিচের সাধারণ-পদ হিসাবটিও ঠিক এই দুর্বল $\beta$-সমতুল্যতাই প্রমাণ করে।} % BN-SRC-312 TARGET-CORRECTION-END

\end{document}
